%! TEX root = ../main.tex

\newpage

% ==============================================================================
% HEADING 1: Tiêu đề chương / mục lớn (Căn trái, 18pt, tự động đánh số 1., 2.,...)
% ==============================================================================
% \section*{Phụ lục A: Các mô hình phân tích}
\addcontentsline{toc}{section}{Phụ lục A: Các mô hình phân tích}

\textit{<Tài liệu này minh họa việc tổ chức các yêu cầu chức năng cho sản phẩm theo các tính năng của hệ thống, các dịch vụ chính được cung cấp bởi sản phẩm.>}

\vspace{0.5cm}

% ==============================================================================
% HEADING 2: Mục con cấp 2 (Căn trái, 15pt, tự động đánh số 1.1, 1.2,...)
% ==============================================================================
\subsection{Tính năng 1 của hệ thống}

<Không nói là "Tính năng 1 của hệ thống" mà viết cụ thể tên chức năng.>

\vspace{0.3cm}

% --- HEADING 3: Mục con cấp 3 (Thụt lề 1.2cm, 13pt, đánh số 1.1.1, 1.1.2,...) ---
\subsubsection{Mô tả và mức ưu tiên}
<Cung cấp một mô tả ngắn gọn về tính năng và chỉ ra mức độ ưu tiên của nó (cao, trung bình, hay thấp)...>

\subsubsection{Tác nhân / Chuỗi đáp ứng}
<Liệt kê các chuỗi hoạt động của người sử dụng và các đáp ứng của hệ thống...>

\subsubsection{Các yêu cầu chức năng}
<Ghi thành từng nhóm các yêu cầu chức năng chi tiết có liên quan tới tính năng này.>

\begin{itemize}
    \item \textbf{REQ-1:} Mô tả yêu cầu chức năng 1.
    \item \textbf{REQ-2:} Mô tả yêu cầu chức năng 2.
\end{itemize}

\vspace{0.5cm}

% ==============================================================================
% HEADING 2 TIẾP THEO
% ==============================================================================
\subsection{Tính năng thứ hai của hệ thống}

Nội dung mô tả tính năng tiếp theo...

% Bảng minh họa (nếu có)
\vspace{0.3cm}

\noindent
\begin{tabularx}{\textwidth}{|p{3.0cm}|X|p{2.5cm}|}
    \hline
    \rowcolor{gray!25}
    \textbf{Mã yêu cầu} & \textbf{Chi tiết} & \textbf{Trạng thái} \\ \hline
    REQ-F01 & Cho phép người dùng đăng nhập hệ thống. & Hoàn thành \\ \hline
    REQ-F02 & Tìm kiếm sách theo từ khóa. & Đang xử lý \\ \hline
\end{tabularx}